\section{Resource-boundary case studies}
\label{sec:s-resource-cases-v2}

The resource coordinate records what was attempted under which envelope and how the benchmark treats the outcome. Table~\ref{tab:s-resource-cases} is the compact index; the five cases below give the reviewer-facing interpretation.

\begingroup
\footnotesize
\setlength{\tabcolsep}{3pt}
\renewcommand{\arraystretch}{1.10}
\begin{longtable}{@{}p{0.21\textwidth} p{0.21\textwidth} p{0.11\textwidth} p{0.15\textwidth} p{0.18\textwidth}@{}}
\caption{Measured resource outcomes and benchmark treatment.}\label{tab:s-resource-cases}\\
\toprule
\textbf{Evaluation cell} & \textbf{Declared envelope / observed event} & \textbf{Outputs} & \textbf{Status} & \textbf{Benchmark treatment} \\
\midrule
\endfirsthead
\caption[]{Measured resource outcomes and benchmark treatment. (continued)}\\
\toprule
\textbf{Evaluation cell} & \textbf{Declared envelope / observed event} & \textbf{Outputs} & \textbf{Status} & \textbf{Benchmark treatment} \\
\midrule
\endhead
\midrule
\multicolumn{5}{r}{\footnotesize Continued on next page}\\
\endfoot
\bottomrule
\endlastfoot
IBM AML HI-Large & safe resource guard & 0 results; 0 predictions & \texttt{SAFE\_\allowbreak{}RESOURCE\_\allowbreak{}BLOCKED} & not in predictive ranking \\
IBM AML LI-Large & safe resource guard & 0 results; 0 predictions & \texttt{SAFE\_\allowbreak{}RESOURCE\_\allowbreak{}BLOCKED} & not in predictive ranking \\
IBM AML HI-Medium GINE h64 & single Tesla T4 CUDA OOM & 0 of 20 planned results; 0 predictions & \texttt{RESOURCE\_\allowbreak{}BLOCKED\_\allowbreak{}T4\_\allowbreak{}CUDA\_\allowbreak{}OOM} & not in predictive ranking \\
IBM AML LI-Medium GINE h64 & single Tesla T4 CUDA OOM & 0 of 20 planned results; 0 predictions & \texttt{RESOURCE\_\allowbreak{}BLOCKED\_\allowbreak{}T4\_\allowbreak{}CUDA\_\allowbreak{}OOM} & not in predictive ranking \\
DGraphFin GAT h64/l2 & Tesla T4 CUDA OOM & 0 of 20 planned results; 0 predictions & \texttt{BLOCKED\_\allowbreak{}T4\_\allowbreak{}OOM} & memory-reduced h32/l1 diagnostic is not a replacement \\
DGraphFin GraphSAGE max-pool rerun & larger GPU required & 0 results; 0 predictions & \texttt{BLOCKED\_\allowbreak{}WAITING\_\allowbreak{}FOR\_\allowbreak{}GPU} & no imported result CSV \\
\end{longtable}
\noindent\footnotesize\emph{Scope and reading.} A blocked cell is unmeasured and excluded from predictive ranks, matched means, and Pareto sets. A reduced diagnostic cannot replace a fixed blocked configuration.\par
\endgroup


\subsection{IBM AML Large guards}

HI-Large and LI-Large stopped at the safe resource guard and produced zero result and prediction files. Evidence exists for the guard and planned configuration; predictive evidence does not exist. The safe wording is that Large lies outside the admitted resource envelope and is unmeasured. It is not permissible to rank Large, infer underperformance, or claim full-10 coverage.

\subsection{Medium GINE T4 out-of-memory}

HI-Medium and LI-Medium GINE h64 each encountered single-T4 CUDA OOM and produced zero of 20 planned outputs. Small GINE performance and the Medium failure record both exist. Medium GINE performance does not. It is excluded from Medium winner sets, matched means, and Pareto sets rather than imputed.

\subsection{Fixed DGraphFin GAT T4 out-of-memory}

The fixed h64/l2 GAT lane has zero of 20 planned results and predictions under the T4 envelope. A memory-reduced h32/l1 diagnostic changes architecture and cannot substitute for the blocked fixed cell. The benchmark may compare their feasibility statuses, not their predictive scores.

\subsection{RB18 and larger-GPU extension status}

The DGraphFin GraphSAGE max-pool extension has preparation commands and status files but no imported result CSV. A larger GPU was requested; no validated runtime record establishes P100, A100, or H100 performance. The safe wording is waiting-for-GPU/unmeasured, not hardware-backed evidence.

\subsection{Memory-reduced exploratory architectures}

Memory-reduced variants can diagnose whether a code path runs under a smaller envelope. They remain exploratory because width, depth, sampling, or features differ from the fixed target. Directory presence, notebook generation, or partial predictions cannot fill a claim scoped to the original architecture.

\subsection{False-promotion protections}

Table~\ref{tab:s-false-promotions-v2} explains how count-only logic could incorrectly promote failed imports, duplicate labels, aliases, diagnostics, and zero-output plans.

\begingroup
\footnotesize
\setlength{\tabcolsep}{3pt}
\renewcommand{\arraystretch}{1.10}
\begin{longtable}{@{}p{0.08\textwidth} p{0.22\textwidth} r r p{0.20\textwidth} p{0.19\textwidth}@{}}
\caption{False-promotion patterns guarded by the artifact layer.}\label{tab:s-false-promotions-v2}\\
\toprule
\textbf{Audit} & \textbf{Family} & \textbf{Result risk} & \textbf{Prediction risk} & \textbf{Correct treatment} & \textbf{Why count-only logic fails} \\
\midrule
\endfirsthead
\caption[]{False-promotion patterns guarded by the artifact layer. (continued)}\\
\toprule
\textbf{Audit} & \textbf{Family} & \textbf{Result risk} & \textbf{Prediction risk} & \textbf{Correct treatment} & \textbf{Why count-only logic fails} \\
\midrule
\endhead
\midrule
\multicolumn{6}{r}{\footnotesize Continued on next page}\\
\endfoot
\bottomrule
\endlastfoot
FP01 & V22 failed/noncomparable imports & 38 & 38 & \texttt{EXCLUDED\_\allowbreak{}INTEGRITY} & missing seed, failed return-code, duplicate/conflicting merge metadata \\
FP02 & V24 RB41 duplicate stress labels & 240 & 240 & \texttt{DEDUPLICATE\_\allowbreak{}TO\_\allowbreak{}120\_\allowbreak{}BASE\_\allowbreak{}CELLS} & stress label is never passed to benchmark harness; all performance metrics repeat \\
FP03 & V24 memory-reduced DGraphFin GAT h32/l1 & 20 & 20 & \texttt{DIAGNOSTIC\_\allowbreak{}ONLY} & different width/depth; canonical lock says not comparable \\
FP04 & V26 legacy P100-named lane & 12 & 12 & \texttt{NORMALIZE\_\allowbreak{}ALIAS\_\allowbreak{}TO\_\allowbreak{}T4/\allowbreak{}CUDA0} & runtime records identify cuda:0 and validation normalizes the alias; path names are not hardware evidence \\
FP05 & IBM AML Large status/runbook artifacts & 0 & 0 & \texttt{RESOURCE\_\allowbreak{}BLOCKED} & canonical lock contains zero results and predictions \\
FP06 & Medium GINE blocked plans & 0 & 0 & \texttt{RESOURCE\_\allowbreak{}BLOCKED} & two T4-OOM rows; zero results and predictions \\
FP07 & RB18 larger-GPU preparation & 0 & 0 & \texttt{BLOCKED\_\allowbreak{}WAITING\_\allowbreak{}FOR\_\allowbreak{}GPU} & validation reports no imported result CSV \\
\end{longtable}
\noindent\footnotesize\emph{Scope and reading.} The 310 result and 310 prediction files at risk are not additional evidence. Zero-output plans remain visible because directories and notebooks can otherwise be mistaken for completed runs.\par
\endgroup


Across metric-bearing artifacts, 310 result and 310 prediction files are at risk of naive promotion. The count includes 240 duplicate V24 labels, 38 failed/noncomparable V22 imports, 20 memory-reduced GAT diagnostics, and 12 legacy hardware-alias files. These remain provenance or validation cases, not additional performance evidence.

\paragraph{General treatment.}
Resource-blocked cells are unmeasured and predictively unordered. A future run on different hardware could produce evidence, but no direction follows from that possibility.
